% =============================================================
%         setting for ACM format paper (Modified from IEEE)
% =============================================================

\usepackage{blkarray}
\usepackage{algpseudocode}
\usepackage{algorithm}
\renewcommand{\algorithmicrequire}{\textbf{Input:}}
\renewcommand{\algorithmicensure}{\textbf{Output:}}
\usepackage{graphicx}

\usepackage[mathcal]{eucal}
\usepackage{mathrsfs}
\usepackage{booktabs}
\usepackage{enumerate}
\usepackage{multirow}
\usepackage[subrefformat=parens,farskip=0pt,justification=centering]{subfig}
\usepackage{color}

\usepackage{comment}
\usepackage{soul}
\soulregister\ref7
\soulregister\pageref7
\usepackage{etoolbox}
\usepackage{url}
\usepackage{nth}
\usepackage{bm}
\usepackage{courier}
\usepackage{balance}
\usepackage{threeparttable}
\usepackage{xcolor,colortbl}
\usepackage{footnote}
\usepackage{listings}
\usepackage{setspace}
\usepackage[inline]{enumitem}
\usepackage{verbatim}

\usepackage{tikz}
\usetikzlibrary{patterns,snakes}
\usetikzlibrary{positioning,calc,fit,decorations.pathmorphing,shapes.geometric, shapes.gates.logic.US, calc}
\usetikzlibrary{arrows,arrows.meta,decorations.markings,shapes,shapes.arrows}
\usetikzlibrary{decorations,decorations.pathreplacing}
\usetikzlibrary{backgrounds}
\usepackage{filecontents}
\usepackage{pgfplots}
\usepackage{pgfplotstable}
\usepackage{scalefnt}
\pgfplotsset{compat=newest}

\usepackage{pifont}
\usepackage{cleveref}
\Crefformat{figure}{Fig.~#2#1#3}
\Crefname{subfigure}{Fig.}{Figs.}
\Crefname{figure}{Fig.}{Figs.}
\Crefformat{table}{TABLE~#2#1#3}

\usepackage[figuresright]{rotating}
\iftrue
\fi

\definecolor{CUHKorange}{RGB}{244,106,18}
\definecolor{CUHKblue}{RGB}{0,111,190}
\definecolor{CUHKgreen}{RGB}{0,127,128}
\definecolor{CUHKred}{RGB}{228,46,36}
\definecolor{CUHKyellow}{RGB}{198,148,34}
\definecolor{CUHKdark}{RGB}{114,44,114}
\definecolor{CUHKmiddle}{RGB}{144,44,144}
\definecolor{CUHKlight}{RGB}{167,44,167} 
\definecolor{CUHKpurple}{RGB}{117,15,109}
\definecolor{CUHKgold}{RGB}{221,163,0}
\definecolor{CUHKribbon}{RGB}{244,223,176}
\definecolor{CUHKblack}{RGB}{34,24,21}

\newcommand{\etal}{\textit{et al.}}
\newcommand{\calH}{\mathcal{H}}
\newcommand{\calL}{\mathcal{L}}
\newcommand{\calN}{\mathcal{N}}
\newcommand{\calO}{\mathcal{O}}
\newcommand{\calP}{\mathcal{P}}
\newcommand{\calQ}{\mathcal{Q}}
\newcommand{\calV}{\mathcal{V}}
\newcommand{\calKL}{\mathcal{KL}}
\newcommand{\norm}[1]{\left\lVert#1\right\rVert}
\newcommand{\tool}[1]{$\mathsf{#1}$}
\renewcommand{\bf}[1]{\textbf{#1}}
\renewcommand{\tt}[1]{\texttt{#1}}
\newcommand{\tensor}[1]{\mathcal{#1}}
\renewcommand{\vec}[1]{\boldsymbol{#1}}
\newcommand{\abs}[1]{\ensuremath{\left\lvert #1\right\rvert}}
\newcommand{\mathbbm}[1]{\text{\usefont{U}{bbm}{m}{n}#1}}
\newcommand{\minisection}[1]{\vspace{.1in}\noindent{\textbf{#1}}.}
\DeclareMathOperator*{\argmin}{argmin}
\DeclareMathOperator*{\argmax}{argmax}
\newcommand\blfootnote[1]{
  \begingroup
  \renewcommand\thefootnote{}\footnote{#1}%
  \addtocounter{footnote}{-1}%
  \endgroup
}
\newcommand*\bcircled[1]{
    \tikz[baseline=(char.base)]{
        \node[shape=circle,fill,inner sep=1pt] (char) {\textcolor{white}{#1}};
    }
}
\newcommand*\wcircled[1]{
    \Large{\textcircled{\scriptsize{\textbf{#1}}}}
}

\usepackage{tcolorbox}
\tcbuselibrary{skins,breakable}
\newenvironment{myblock}[1]{%
    \tcolorbox[beamer,%
    noparskip,breakable,
    colback=white,colframe=CUHKmiddle,%
    colbacklower=white,%
    title=#1]}%
    {\endtcolorbox}
\newenvironment{commentblock}[1]{%
    \tcolorbox[beamer,%
    noparskip,breakable,
    colback=white,colframe=CUHKmiddle,%
    colbacklower=white,%
    title=Comment #1]}%
    {\endtcolorbox}

\newtheorem{myproblem}{\textbf{Problem}}
\newtheorem{mydefinition}{\textbf{Definition}}
\newtheorem{mytheorem}{\textbf{Theorem}}
\newtheorem{mycorollary}{\textbf{Corollary}}
\newtheorem{mylemma}{\textbf{Lemma}}
\newtheorem{myclaim}{\textbf{Claim}}
\newtheorem{myapplication}{\textbf{Application}}
\newtheorem{myexample}{\textbf{Example}}
\newtheorem{myconjecture}{\textbf{Conjecture}}
\newtheorem{myassumption}{Assumption}
\newtheorem{myprop}{Property}

\crefname{mytheorem}{Theorem}{Theorems}
\crefname{mylemma}{Lemma}{Lemmas}
\crefname{myclaim}{Claim}{Claims}
\crefname{myproperty}{Property}{Properties}
\crefname{mycorollary}{Corollary}{Corollaries}

\algrenewcommand\textproc{\texttt}
\newcommand{\tabincell}[2]{\begin{tabular}{@{}#1@{}}#2\end{tabular}}

\makeatletter
\let\OldStatex\Statex
\renewcommand{\Statex}[1][3]{%
  \setlength\@tempdima{\algorithmicindent}%
  \OldStatex\hskip\dimexpr#1\@tempdima\relax
}
\makeatother

\RequirePackage[normalem]{ulem} 
\RequirePackage{color}\definecolor{RED}{rgb}{1,0,0}\definecolor{BLUE}{rgb}{0,0,1} 
\providecommand{\DIFadd}[1]{{\protect\color{blue}{#1}}}                           
\providecommand{\DIFdel}[1]{{\protect\color{red}\sout{#1}}}                       
\providecommand{\DIFaddbegin}{}                                                   
\providecommand{\DIFaddend}{}                                                     
\providecommand{\DIFdelbegin}{}                                                   
\providecommand{\DIFdelend}{}                                                     
\providecommand{\DIFaddFL}[1]{\DIFadd{#1}}                                        
\providecommand{\DIFdelFL}[1]{\DIFdel{#1}}                                        
\providecommand{\DIFaddbeginFL}{}                                                 
\providecommand{\DIFaddendFL}{}                                                   
\providecommand{\DIFdelbeginFL}{}                                                 
\providecommand{\DIFdelendFL}{}                                                   

\newcommand{\todo}[1]{\textcolor{red}{[TODO: #1]}}
\newcommand{\revise}[1]{\DIFadd{#1}}